\documentclass[10pt,aspectratio=169]{beamer}
% Preámbulo modular para presentaciones Beamer (Modelación Estadística)
% Diseño y paleta institucional del Tecnológico de Monterrey

\documentclass[aspectratio=169,xcolor={dvipsnames,table}]{beamer}

\usepackage[utf8]{inputenc}
\usepackage[spanish,mexico]{babel}
\usepackage{amsmath,amssymb,amsfonts,mathtools}
\usepackage{graphicx}
\usepackage{listings}
\usepackage{booktabs}
\usepackage{tikz}
\usetikzlibrary{matrix,arrows,decorations.pathmorphing,shapes.geometric,positioning}

% Paleta Institucional Tecnológico de Monterrey
\definecolor{TecAzul}{HTML}{0039A6}       % Azul TEC: color primario institucional
\definecolor{TecAzulOscuro}{HTML}{1A2E51} % Azul marino oscuro
\definecolor{TecRojo}{HTML}{EC2661}       % Rojo de acento (paleta miTec)
\definecolor{TecGris}{HTML}{646464}       % Gris neutro para comentarios y subtítulos
\definecolor{TecFondoBloque}{HTML}{F4F6F9}% Fondo suave para bloques

% Configuración de Tema Beamer
\usetheme{Boadilla}
\usecolortheme[named=TecAzulOscuro]{structure}
\setbeamercolor{alerted text}{fg=TecRojo}
\setbeamercolor{palette primary}{bg=TecAzulOscuro,fg=white}
\setbeamercolor{palette secondary}{bg=TecAzul,fg=white}
\setbeamercolor{palette tertiary}{bg=TecAzulOscuro!80!black,fg=white}
\setbeamercolor{title}{fg=TecAzulOscuro,bg=TecFondoBloque}
\setbeamercolor{frametitle}{fg=TecAzulOscuro,bg=TecFondoBloque}

% Bloques personalizados elegantes
\setbeamercolor{block title}{bg=TecAzulOscuro,fg=white}
\setbeamercolor{block body}{bg=TecFondoBloque,fg=black}
\setbeamercolor{block title alerted}{bg=TecRojo,fg=white}
\setbeamercolor{block body alerted}{bg=TecRojo!8,fg=black}
\setbeamercolor{block title example}{bg=TecAzul,fg=white}
\setbeamercolor{block body example}{bg=TecAzul!8,fg=black}

% Redondear bordes y añadir sombra sutil
\setbeamertemplate{blocks}[rounded][shadow=true]
\setbeamertemplate{navigation symbols}{}

% Pie de página limpio con número de diapositiva
\setbeamertemplate{footline}{
  \leavevmode%
  \hbox{%
  \begin{beamercolorbox}[wd=.7\paperwidth,ht=2.25ex,dp=1ex,left]{author in head/foot}%
    \hspace*{2ex}\insertshortauthor~~(\insertshortinstitute) \hfill \insertshorttitle
  \end{beamercolorbox}%
  \begin{beamercolorbox}[wd=.3\paperwidth,ht=2.25ex,dp=1ex,right]{date in head/foot}%
    \insertshortdate{}\hspace*{2em}
    \insertframenumber{} / \inserttotalframenumber\hspace*{2ex}
  \end{beamercolorbox}}%
  \vskip0pt%
}

% Configuración de Listings de Python para visualización pedagógica de código
\lstset{
    language=Python,
    basicstyle=\ttfamily\footnotesize,
    keywordstyle=\color{TecAzulOscuro}\bfseries,
    stringstyle=\color{TecRojo},
    commentstyle=\color{TecGris}\itshape,
    showstringspaces=false,
    breaklines=true,
    frame=lines,
    rulecolor=\color{TecAzulOscuro!30},
    backgroundcolor=\color{TecFondoBloque},
    aboveskip=0.5em,
    belowskip=0.5em,
    literate={á}{{\'a}}1 {é}{{\'e}}1 {í}{{\'i}}1 {ó}{{\'o}}1 {ú}{{\'u}}1
             {Á}{{\'A}}1 {É}{{\'E}}1 {Í}{{\'I}}1 {Ó}{{\'O}}1 {Ú}{{\'U}}1
             {ñ}{{\~n}}1 {Ñ}{{\~N}}1 {¿}{{?`}}1 {¡}{{!`}}1
}

% Metadatos institucionales por defecto
\title[Modelación Estadística]{Modelación Estadística para Ciencia de Datos e Ingeniería}
\author[J. Castillo Colmenares]{Juliho David Castillo Colmenares}
\institute[Tec de Monterrey]{Tecnológico de Monterrey \\ Escuela de Ingeniería y Ciencias}
\date{\today}

% Comandos y macros estadísticas compartidas para las presentaciones Beamer (Metropolis)

% Conjuntos numéricos básicos
\newcommand{\R}{\mathbb{R}}
\newcommand{\N}{\mathbb{N}}
\newcommand{\Q}{\mathbb{Q}}
\newcommand{\Z}{\mathbb{Z}}
\newcommand{\set}[1]{\left\{ #1 \right\}}
\newcommand{\abs}[1]{\left|#1\right|}
\newcommand{\norm}[1]{\left\| #1 \right\|}

% Comandos estadísticos y probabilísticos del curso
\newcommand{\Var}{\operatorname{Var}}
\newcommand{\cov}{\operatorname{Cov}}
\newcommand{\corr}{\rho}
\newcommand{\comb}[2]{\begin{pmatrix} #1 \\ #2 \end{pmatrix}}
\newcommand{\s}{\sigma}
\newcommand{\particion}{\mathfrak{P}}
\newcommand{\card}[1]{\#\left( #1 \right)}
\newcommand{\seq}[3]{\set{#1}_{#2}^{#3}}
\newcommand{\E}{\mathbb{E}}
\newcommand{\Prob}{\mathbb{P}}

% Entornos pedagógicos y cajas en español académico natural
\newenvironment{discusion}{%
  \begin{alertblock}{Pregunta de Discusión}%
}{%
  \end{alertblock}%
}

\newenvironment{reflexion}{%
  \begin{alertblock}{Reflexión para el Aula}%
}{%
  \end{alertblock}%
}

\newenvironment{paradoja}{%
  \begin{alertblock}{Paradoja de Exploración}%
}{%
  \end{alertblock}%
}

\newenvironment{motivacion}{%
  \begin{exampleblock}{Motivación: Ciencia de Datos en la Práctica}%
}{%
  \end{exampleblock}%
}

\newenvironment{retoclase}{%
  \begin{block}{Reto Rápido en Equipo}%
}{%
  \end{block}%
}


\title[Tamaño de una muestra]{Sección 6.7: Estimación del tamaño de una muestra}
\subtitle{Confianza, margen de error y variabilidad previa}
\author[J. Castillo Colmenares]{Juliho Castillo Colmenares}
\institute{Tecnológico de Monterrey}
\date{}

\begin{document}
\begin{frame}[plain]\vspace{-0.8cm}\titlepage\end{frame}

\begin{frame}{Hoja de ruta --- capítulo 6}
  \footnotesize
  \begin{itemize}\itemsep=0.08cm
    \item \textbf{06.02} Intervalos de confianza
    \item \textbf{06.04--06.05} Errores estándar y proporciones
    \item \textbf{06.07} Tamaño de muestra \emph{(hoy)}
    \item \textbf{07} Pruebas de hipótesis y potencia
  \end{itemize}
  \begin{block}{Objetivo didáctico}
    Determinar $n$ antes de recolectar datos fijando confianza, margen de error y
    una estimación previa de la variabilidad.
  \end{block}
\end{frame}

\begin{frame}{Motivación: ¿cuántos datos?}
  \footnotesize
  Una muestra muy pequeña puede carecer de poder; una demasiado grande
  desperdicia recursos financieros, tiempo e infraestructura.
  \begin{alertblock}{Pregunta guía}
    ¿Qué tamaño garantiza que el margen de error no exceda $E$?
  \end{alertblock}
\end{frame}

\begin{frame}{Tres decisiones previas}
  \footnotesize
  Antes de recolectar datos se especifican:
  \begin{enumerate}\itemsep=0.12cm
    \item confianza nominal $(1-\alpha)$;
    \item margen máximo tolerable $E$;
    \item dispersión previa $\sigma$ o proporción esperada $p^*$.
  \end{enumerate}
\end{frame}

\begin{frame}{Tamaño para una media}
  \footnotesize
  Como $E=Z_{\alpha/2}\sigma/\sqrt n$,
  \[
    n=\left(\frac{Z_{\alpha/2}\sigma}{E}\right)^2.
  \]
  Si $\sigma$ es desconocida, puede estimarse con un piloto o con la regla del
  rango $\sigma\approx\text{rango}/4$.
\end{frame}

\begin{frame}{Tamaño para una proporción}
  \footnotesize
  Como $E=Z_{\alpha/2}\sqrt{p(1-p)/n}$,
  \[
    n=\left(\frac{Z_{\alpha/2}}{E}\right)^2p^*(1-p^*).
  \]
  El resultado se redondea siempre hacia arriba.
\end{frame}

\begin{frame}{Caso conservador y población finita}
  \footnotesize
  Sin información previa, $p^*=0.5$ maximiza $p(1-p)$:
  \[
    n_{\max}=\frac{Z_{\alpha/2}^2}{4E^2}.
  \]
  Si $n/N>0.05$, se ajusta por población finita:
  \[
    n_a=\frac{n}{1+(n-1)/N}.
  \]
\end{frame}

\begin{frame}{Monotonicidad del tamaño requerido}
  \footnotesize
  El tamaño aumenta con $Z_{\alpha/2}$ y la variabilidad, y disminuye cuando se
  tolera un margen $E$ mayor.
  \begin{alertblock}{Lectura}
    Las fórmulas hacen explícitos los intercambios entre precisión y recursos.
  \end{alertblock}
\end{frame}

\begin{frame}{Puente computacional 1/4: efecto de confianza}
  \footnotesize
  Un nivel de confianza mayor aumenta $Z_{\alpha/2}$ y, por tanto, el tamaño
  requerido.
  \begin{block}{Regla}
    Más confianza o menor $E$ implica más observaciones.
  \end{block}
\end{frame}

\begin{frame}{Puente computacional 2/4: efecto del margen}
  \footnotesize
  \[
    n\propto\frac1{E^2}.
  \]
  Reducir el margen a la mitad cuadruplica aproximadamente el tamaño de muestra.
\end{frame}

\begin{frame}{Puente computacional 3/4: variabilidad}
  \footnotesize
  \begin{columns}[T]
    \column{0.48\textwidth}
      Media: $n\propto\sigma^2$.
    \column{0.48\textwidth}
      Proporción: $n\propto p^*(1-p^*)$.
  \end{columns}
  \vspace{0.2cm}
  La información previa reduce el costo cuando es fiable.
\end{frame}

\begin{frame}{Puente computacional 4/4: redondeo y FPC}
  \footnotesize
  \begin{enumerate}\itemsep=0.12cm
    \item calcular la fórmula continua;
    \item redondear hacia arriba;
    \item aplicar FPC si la fracción muestral supera $5\%$.
  \end{enumerate}
\end{frame}

\begin{frame}{Aplicación de lectura 1/4 --- media}
  \footnotesize
  Se conoce una estimación previa $\sigma$, se fija $E$ y se selecciona un
  nivel de confianza.
  \begin{block}{Planteamiento}
    Despejar $n$ del margen de error de la media.
  \end{block}
\end{frame}

\begin{frame}{Aplicación de lectura 1/4 --- solución}
  \footnotesize
  \[
    n=\left(\frac{Z_{\alpha/2}\sigma}{E}\right)^2,
    \qquad n_{\text{final}}=\lceil n\rceil.
  \]
  El redondeo hacia arriba garantiza $E_{\text{final}}\le E$.
\end{frame}

\begin{frame}{Aplicación de lectura 2/4 --- proporción con $p^*$}
  \footnotesize
  Un estudio de respuestas binarias dispone de una proporción esperada $p^*$.
  \begin{block}{Planteamiento}
    Usar la variabilidad anticipada para calcular el tamaño requerido.
  \end{block}
\end{frame}

\begin{frame}{Aplicación de lectura 2/4 --- solución}
  \footnotesize
  \[
    n=\left(\frac{Z_{\alpha/2}}{E}\right)^2p^*(1-p^*).
  \]
  Una proporción previa alejada de $0.5$ reduce la varianza esperada.
\end{frame}

\begin{frame}{Aplicación de lectura 3/4 --- caso conservador}
  \footnotesize
  Si no existe información previa sobre $p$, debe protegerse el peor caso.
  \begin{block}{Planteamiento}
    Maximizar $p(1-p)$ en $[0,1]$.
  \end{block}
\end{frame}

\begin{frame}{Aplicación de lectura 3/4 --- solución}
  \footnotesize
  \[
    p^*=0.5,\qquad p^*(1-p^*)=0.25,
    \qquad n_{\max}=\frac{Z_{\alpha/2}^2}{4E^2}.
  \]
  Este tamaño no subestima la variabilidad posible de una proporción.
\end{frame}

\begin{frame}{Aplicación de lectura 4/4 --- población finita}
  \footnotesize
  Una población de tamaño $N$ puede ser pequeña frente al tamaño calculado.
  \begin{block}{Planteamiento}
    Revisar si $n/N>0.05$ y ajustar el tamaño.
  \end{block}
\end{frame}

\begin{frame}{Aplicación de lectura 4/4 --- solución}
  \footnotesize
  \[
    n_a=\frac{n}{1+(n-1)/N}.
  \]
  El ajuste reduce observaciones sin sacrificar el margen de error estipulado.
\end{frame}

\begin{frame}{Síntesis de la sección}
  \footnotesize
  \begin{itemize}\itemsep=0.12cm
    \item Confianza, margen y variabilidad determinan $n$.
    \item El margen entra al cuadrado en el denominador.
    \item $p=0.5$ es el caso conservador para proporciones.
    \item La FPC corrige poblaciones finitas.
  \end{itemize}
\end{frame}

\begin{frame}{Transición curricular}
  \footnotesize
  \begin{block}{Lo que queda establecido}
    El tamaño muestral se decide antes de observar datos, con supuestos
    explícitos y redondeo seguro.
  \end{block}
  \begin{alertblock}{Siguiente unidad}
    Las pruebas de hipótesis usarán estos tamaños para controlar significación,
    potencia y precisión.
  \end{alertblock}
\end{frame}
\end{document}
